%======================================================== 6 DANH SÁCH
\begin{frame}[label=danhsach]{4.\ Danh sách bài và bộ lọc}
\begin{tikzpicture}[x=1mm,y=1mm]
\browser{0}{0}{92}{62}
\node[anchor=west,font=\fsz{7}{8}\bfseries] (h) at (2,-7) {Đề ra kỳ này};
\node[anchor=west,font=\fsz{4.6}{5.4},text=ink!60] (me) at (2,-10.5) {ban@example.com · NCB};
\node[sel,anchor=west] (f1) at (27,-10.5) {Mọi chủ đề {\symf ⌄}};
\node[sel,anchor=west] (f2) at (42,-10.5) {Mọi mức {\symf ⌄}};
\node[sel,anchor=west] (f3) at (54,-10.5) {Mọi trạng thái {\symf ⌄}};
\node[sel,anchor=west,draw=QT] (f4) at (72,-10.5) {Vai trò thật {\symf ⌄}};
\draw[ink!15] (0,-13.5) -- (92,-13.5);
% card 1
\draw[ink!20,fill=white,rounded corners=1pt] (6,-16) rectangle (86,-34);
\node[anchor=west,font=\fsz{5.6}{6.4}\ttfamily\bfseries,text=PB!70!black] (code) at (8,-19.5) {2026-07-06a};
\node[pchip,anchor=west] (c1) at (26,-19.5) {ĐS};
\node[pchip,anchor=west] (c2) at (31,-19.5) {Mức B};
\node[pchip,anchor=west] (c3) at (40,-19.5) {SL-Fail};
\node[wchip,anchor=west] (c4) at (48,-19.5) {cần kiểm tra 2};
\node[wchip,anchor=west] (c5) at (61.5,-19.5) {xung đột mở 1};
\node[anchor=east,font=\fsz{4.8}{5.6}\itshape,text=ink!60] (au) at (84.5,-19.5) {Tác Giả};
\node[anchor=north west,font=\fsz{5.4}{6.6},text width=74mm] at (8,-22.5) {Cho $a,b,c$ là các số thực dương thoả mãn $a+b+c=3$. Chứng minh rằng \dots};
\draw[ink!20,fill=white,rounded corners=1pt] (6,-36) rectangle (86,-48);
\node[anchor=west,font=\fsz{5.6}{6.4}\ttfamily\bfseries,text=PB!70!black] at (8,-39.5) {SL09-1047};
\node[pchip,anchor=west] at (26,-39.5) {ĐS}; \node[pchip,anchor=west] at (31,-39.5) {Mức B}; \node[pchip,anchor=west] at (40,-39.5) {SL-Fail};
\fill[ink!15] (8,-43) rectangle (80,-43.8); \fill[ink!15] (8,-45.5) rectangle (60,-46.3);
\draw[ink!20,fill=white,rounded corners=1pt] (6,-50) rectangle (86,-60);
\fill[ink!15] (8,-55) rectangle (80,-55.8);
% callouts
\begin{scope}[every node/.style={callout=ink,anchor=west,text width=44mm}]
\node (k1) at (99,-4) {\textbf{Email · vai trò} của bạn. Quản trị có thêm ô \emph{Vai trò thật} để xem như vai trò khác.};
\node (k2) at (99,-14) {\textbf{Bộ lọc}: chủ đề (ĐS, SH, HH, TH), mức (A, B), trạng thái. Kết hợp được, ví dụ HH + A + SL-OK.};
\node (k3) at (99,-25.5) {\textbf{Mã bài} — bấm vào để mở trang bài.};
\node (k4) at (99,-34) {\textbf{Thuộc tính}: chủ đề · mức · trạng thái (mục 5).};
\node[callout=warn] (k5) at (99,-43.5) {\textbf{Việc còn mở}: số mục cần kiểm tra, số xung đột chưa giải quyết. Không hiện với phản biện.};
\node (k6) at (99,-54) {\textbf{Tác giả}: ẩn với phản biện (chấm ẩn danh).};
\end{scope}

\draw[cline=ink!50] (k1.west) -- (me.north east);
\draw[cline=ink!50] (k2.west) -- (f3.north east);
\draw[cline=ink!50] (k3.west) -- (code.south east);
\draw[cline=ink!50] (k4.west) -- (c3.south);
\draw[cline=warn!70] (k5.west) -- (c5.south);
\draw[cline=ink!50] (k6.west) -- (au.south);
\end{tikzpicture}
\end{frame}

%======================================================== 7 MÃ BÀI
\begin{frame}[label=mabai]{5.\ Mã bài và thuộc tính}
\begin{tikzpicture}[x=1mm,y=1mm]
\node[seg=TAGG] (c0) at (0,-8) {2026-06-10a};
\node[seg=tHH,right=4mm of c0] (c1) {HH};
\node[seg=PT,right=4mm of c1] (c2) {A};
\node[seg=PB,right=4mm of c2] (c3) {SL-OK};
\node[seg=BTK,right=4mm of c3] (c4) {Pi 11/2026 · P1045};
\draw[decorate,decoration={brace,amplitude=2mm},line width=0.7pt,ink!60] ([yshift=1mm]c0.north west) -- ([yshift=1mm]c0.north east) node[midway,above=2.5mm,font=\fsz{6}{7}\bfseries]{MÃ BÀI — không bao giờ đổi};
\draw[decorate,decoration={brace,amplitude=2mm},line width=0.7pt,ink!60] ([yshift=1mm]c1.north west) -- ([yshift=1mm]c4.north east) node[midway,above=2.5mm,font=\fsz{6}{7}\bfseries]{THUỘC TÍNH — thay đổi theo vòng đời của bài};
\node[anchor=north,font=\fsz{5.4}{6.4},text=ink!70,align=center] at (c0.south) {năm-tháng nhận · số hồ sơ\\a, b, c… = bài thứ mấy trong hồ sơ};
\node[anchor=north,font=\fsz{5.4}{6.4},text=tHH!70!black] at (c1.south) {chủ đề};
\node[anchor=north,font=\fsz{5.4}{6.4},text=PT] at (c2.south) {mức};
\node[anchor=north,font=\fsz{5.4}{6.4},text=PB] at (c3.south) {trạng thái};
\node[anchor=north,font=\fsz{5.4}{6.4},text=BTK] at (c4.south) {số tạp chí · số in (khi đã đăng)};
\draw[ink!25,fill=soft,rounded corners=2pt] (0,-24) rectangle (88,-62);
\begin{scope}[every node/.style={anchor=west,font=\fsz{6.2}{7.4}}]
\node at (2,-28) {\textbf{Chủ đề}}; \node at (20,-28) {ĐS đại số · SH số học · HH hình học · TH tổ hợp};
\node at (2,-33) {\textbf{Mức}}; \node at (20,-33) {A · B};
\node at (2,-38) {\textbf{Trạng thái}}; \node at (20,-38) {\textbf{Mới} — đã nhập, chưa vào shortlist};
\node at (20,-42.5) {\textbf{SL} — trong shortlist, đang phản biện};
\node at (20,-47) {\textbf{SL-OK} — đạt, chờ đăng (dự trữ)};
\node at (20,-51.5) {\textbf{SL-Fail} — chưa đạt ở lần xét này};
\node at (20,-56) {\textbf{PL} — đã đăng};
\end{scope}
\node[anchor=north west,text width=50mm,font=\fsz{6.4}{8}] at (93,-24) {%
\textbf{Mã tạm.} Bài chưa tìm được hồ sơ gốc dùng mã tạm, ví dụ \texttt{SL09-1042} (bảng chọn bài tháng 9, số 1042). Khi tìm được hồ sơ, bài nhận mã chính thức; mã tạm vẫn được ghi lại.\\[3mm]
\textbf{Bài luyện tập} \texttt{THU-01…} là bài bịa để thử hệ thống; chỉ người được giao và Quản trị thấy.\\[3mm]
\textbf{Số in} (P1045…) chỉ có khi khoá kỳ. Các danh sách cũ đánh số khác nhau cho cùng một bài; mã bài chấm dứt việc đó.};
\end{tikzpicture}
\end{frame}

%======================================================== 8 TRANG BÀI
\begin{frame}[label=trangbai]{6.\ Trang một bài}
\begin{tikzpicture}[x=1mm,y=1mm]
\browser{0}{0}{92}{64}
\node[anchor=west,font=\fsz{5}{6},text=NCB] (back) at (3,-6.5) {← Danh sách};
\node[anchor=west,font=\fsz{5.6}{6.4}\ttfamily\bfseries] at (3,-10.5) {2026-09-11a};
\node[pchip,anchor=west] at (21,-10.5) {ĐS}; \node[pchip,anchor=west] at (26,-10.5) {Mức B}; \node[pchip,anchor=west] at (35,-10.5) {Mới};
\node[anchor=west,font=\fsz{6}{7}\bfseries] (de) at (3,-15.5) {Đề bài};
\fill[ink!15] (3,-18.5) rectangle (86,-19.3); \fill[ink!15] (3,-21) rectangle (70,-21.8);
\node[anchor=west,font=\fsz{6}{7}\bfseries] (lg) at (3,-26) {Lời giải};
\foreach \y in {-29,-31.5,-34,-36.5}{\fill[ink!15] (3,\y) rectangle (84,\y-0.8);}
\node[anchor=west,font=\fsz{4.8}{5.6},text=warn] (cb) at (3,-40) {Cảnh báo hiển thị: lệnh chưa hỗ trợ …};
\node[anchor=west,font=\fsz{5}{6}\bfseries] (ng) at (3,-44.5) {{\symf ▸} Nguồn \& chỉnh sửa};
\node[anchor=west,font=\fsz{6}{7}\bfseries] (tl) at (3,-49.5) {Thảo luận};
\node[anchor=west,font=\fsz{4.8}{5.6}] at (3,-53) {\textbf{pb1@example.com}\ \textcolor{ink!50}{2026-10-12}\ \ Bước cuối nên viết rõ \dots};
\draw[ink!30,fill=white] (3,-55.5) rectangle (70,-61.5);
\node[anchor=west,font=\fsz{4.4}{5},text=ink!45] at (4,-57) {Nhận xét (có thể viết công thức \$…\$)};
\node[btn,anchor=west] (gui) at (72,-58.5) {Gửi nhận xét};
\begin{scope}[every node/.style={callout=ink,anchor=west,text width=44mm}]
\node (k1) at (99,-6) {Về danh sách. (Không dùng nút tải lại của trình duyệt.)};
\node (k2) at (99,-17) {\textbf{Đề bài, lời giải}: bản đã biên tập, công thức hiện như khi in.};
\node[callout=warn] (k3) at (99,-29) {Dòng đỏ này chỉ hiện khi văn bản có lệnh \LaTeX\ hệ thống chưa hiển thị được — báo NCB.};
\node (k4) at (99,-40.5) {\textbf{Nguồn \& chỉnh sửa}: bấm để mở (giải thích ở phần NCB). Phản biện không có mục này.};
\node (k5) at (99,-53) {\textbf{Thảo luận}: mọi người xem được bài đều thấy và trả lời được (mục 7).};
\end{scope}
\draw[cline=ink!50] (k1.west) -- (back.east);
\draw[cline=ink!50] (k2.west) -- (de.east);
\draw[cline=warn!70] (k3.west) -- (cb.east);
\draw[cline=ink!50] (k4.west) -- (ng.east);
\draw[cline=ink!50] (k5.west) -- (tl.east);
\end{tikzpicture}
\end{frame}

%======================================================== 9 NHẬN XÉT
\begin{frame}[label=nhanxet]{7.\ Viết nhận xét có công thức}
\begin{tikzpicture}[x=1mm,y=1mm]
\node[anchor=north west,font=\fsz{6.4}{7.6}\bfseries] at (0,0) {Bạn gõ};
\node[anchor=north west,font=\fsz{6.4}{7.6}\bfseries] at (76,0) {Mọi người thấy};
\foreach \src/\out [count=\i] in {
  {Bước 2: \$a+b\textbackslash ge 2\textbackslash sqrt\{ab\}\$ đúng vì \dots}/{Bước 2: $a+b\ge 2\sqrt{ab}$ đúng vì \dots},
  {\$\$\textbackslash frac\{x\}\{y\}+\textbackslash frac\{y\}\{x\}\textbackslash ge 2\$\$}/{\hspace*{18mm}$\displaystyle\frac{x}{y}+\frac{y}{x}\ge 2$},
  {Đáp số đúng là \$n=5\$, không phải \$14\$.}/{Đáp số đúng là $n=5$, không phải $14$.},
  {Mã HTML như <b>đậm</b> không có tác dụng.}/{Mã HTML như \texttt{<b>đậm</b>} không có tác dụng.},
  {\textbackslash textbf\{Ý chính\}: bài nên ở mức A.}/{\textbf{Ý chính}: bài nên ở mức A.}}{
  \node[draw=ink!25,fill=white,rounded corners=1pt,anchor=north west,text width=68mm,minimum height=8mm,font=\fsz{5.8}{7}\ttfamily] at (0,{-6-(\i-1)*10}) {\src};
  \node[font=\fsz{8}{8},text=ink!40] at (72,{-10-(\i-1)*10}) {→};
  \node[draw=PB!40,fill=PB!4,rounded corners=1pt,anchor=north west,text width=66mm,minimum height=8mm,font=\fsz{6.4}{7.6}] at (76,{-6-(\i-1)*10}) {\out};}
\node[anchor=north west,text width=144mm,font=\fsz{6.4}{8}] at (0,-57) {%
\textbf{Công thức}: \texttt{\$…\$} trong dòng, \texttt{\$\$…\$\$} riêng một dòng — giống \LaTeX. Các lệnh riêng của Pi (\texttt{\textbackslash dtr}, \texttt{\textbackslash vto}…) dùng được.
\textbf{Mã HTML không có tác dụng} (hiện ra như chữ). Nhận xét \textbf{chưa sửa/xoá được} sau khi gửi — đọc lại trước khi bấm \emph{Gửi nhận xét}.
Một dòng trống tách đoạn. Ban biên tập thấy email người viết; phản biện thấy tên ẩn danh (\emph{Bạn}, \emph{Phản biện 1, 2…}, \emph{Ban biên tập}).};
\end{tikzpicture}
\end{frame}

%======================================================== 10 AI THẤY GÌ
\begin{frame}[label=aithay]{8.\ Ai thấy gì}
\vskip-1mm
{\fsz{6.6}{8.2}
\setlength{\tabcolsep}{4pt}\renewcommand{\arraystretch}{1.35}
\begin{tabular}{@{}l*{7}{c}@{}}
 & \rtag{TBT} & \rtag{PT} & \rtag{NCB} & \rtag{VP} & \rtag{BTK} & \rtag{PB} & \rtag{QT}\\\hline
Danh sách bài & mọi bài & mọi bài & mọi bài & mọi bài & mọi bài & \textbf{bài được giao} & mọi bài\\
Tên tác giả & \yes & \yes & \yes & \yes & \yes & \no\ (ẩn danh) & \yes\\
Liên hệ tác giả & \yes & \yes & \no & \yes & \no & \no & \yes\\
Nguồn \& chỉnh sửa, xung đột & \yes & \yes & \yes & \yes & \yes & \no & \yes\\
Lịch sử sửa đề, lời giải & \yes & \yes & \yes & \yes & \yes & \no & \yes\\
Thảo luận (xem, viết) & \yes & \yes & \yes & \yes & \yes & bài được giao & \yes\\\hline
Sửa đề, lời giải & & \yes & \yes & & & & \yes\\
Mục cần kiểm tra, xung đột, sửa đổi & \yes & \yes & \yes & & & & \yes\\
Đổi trạng thái bài & \yes & \yes & & & & & \yes\\
Phiếu phản biện (xem) & \yes & \yes & \yes & & & của mình & \yes\\
Kỳ phản biện: mở, giao bài, thư mời & & \yes & & & & & \yes\\
Thêm người dùng, giao vai trò & & & & & & & \yes\\\hline
\end{tabular}}
\vskip3mm
{\fsz{6.2}{7.6}\color{ink!75}
\textbf{Phản biện} chỉ thấy bài được giao trong \textbf{kỳ đang mở}; hết kỳ, bài không còn trong danh sách của họ.
Phản biện không thấy tên tác giả, tên tệp, nguồn, xung đột — để chấm ẩn danh.
Mỗi lần mở một bài đều được ghi lại (ai, bài nào, khi nào).
Xung đột \emph{mức}, \emph{tác giả}, \emph{trùng bài}: chỉ TBT ghi cách giải quyết; người khác chuyển sang „chờ TBT".}
\end{frame}
